\section{One original leading function and the cancellation of local parameters}

Let $i:C\hookrightarrow Y$ be a genuine closed immersion of integral schemes.
Put $\eta=i(\eta_C)$, $D=\mathcal O_{Y,\eta}$, $K=Y.\mathrm{functionField}$,
and $F=C.\mathrm{functionField}$. The homomorphism
\[
 \lambda:D\longrightarrow F
\]
is the ORIGINAL stalk map $i^\#_{\eta_C}$. At an original $x\in C$ put
$S_x=\mathcal O_{Y,i(x)}$. An original component equation $\pi_x\in S_x$
is prime and satisfies $\ker(i_x^\#)=(\pi_x)$. Write $T_x=(S_x)_{(\pi_x)}$
and retain the actual original specialization map $j_x:S_x\to D$.
The existing construction gives $E_x:T_x\simeq D$, with
$E_x(s/1)=j_x(s)$ for EVERY $s\in S_x$.

For the residual-order statements assume $S_x$ Noetherian and choose the
original second parameter $z_x$ with $\pi_x\nmid z_x$ and
$\mathfrak m_{S_x}=(\pi_x,z_x)$. The existing original leading-field
isomorphism is denoted
$R_x:\kappa(T_x)\simeq F$. Its target is the SAME original function
field of $C$, rather than the residue field of the closed point $x$.
Let $J_x:\operatorname{Frac}S_x\simeq K$ be the canonical original
fraction-field isomorphism and let $\iota_x:T_x\to\operatorname{Frac}S_x$
be the existing original vertical fraction inclusion.

\begin{lemma}[Both original whole-localization squares]
\label{lem:gauss-leading-original-localization-squares}
\uses{thm:gauss-common-original-component-stalk,def:central-curve-original-leading-field,def:gauss-surface-original-vertical-fraction-map,lem:gauss-surface-original-vertical-fraction-identities}
\lean{EnvelopeProperty.GaussLeadingResidueNaturality.componentStalk_residual_apply,EnvelopeProperty.GaussLeadingResidueNaturality.componentStalk_fraction_apply}
\leanok
For EVERY $w\in T_x$, including zero,
\[
 R_x(\operatorname{res}_{T_x}w)=\lambda(E_xw),\qquad
 J_x(\iota_xw)=\iota_D(E_xw),
\]
where $\iota_D:D\to K$ is the original stalk-to-function-field map.
Neither displayed identity is an input.
\end{lemma}
\begin{proof}
Both sides of the first equality define ring homomorphisms $T_x\to F$.
On an original coefficient $s\in S_x$, the existing leading-field
coefficient formula identifies the left side with the original image of
$i_x^\#s$ in $F$. The actual naturality of stalk maps for
$\eta_C\leadsto x$ identifies this with $\lambda(j_xs)$. The existing
coefficient formula for $E_x$ gives the right side. Uniqueness of a ring
homomorphism out of the localization of $S_x$ therefore proves equality
on ALL of $T_x$.

For the second equality both sides are ring homomorphisms $T_x\to K$.
On $s\in S_x$, the original vertical inclusion and the canonical
fraction-field isomorphism give its original image in $K$. Composition
of the actual structure-sheaf specialization maps shows that the same
image is $\iota_D(j_xs)$. Apply the same localization uniqueness. Thus
the proof retains the entire original coefficient action and not only
the values of selected parameters or units.
\end{proof}

\begin{lemma}[The same exact vertical order on the whole original field]
\label{lem:gauss-leading-original-vertical-order}
\uses{lem:gauss-leading-original-localization-squares,thm:gauss-common-original-component-stalk,lem:central-curve-fraction-order-ring-equiv,def:gauss-surface-same-field-native-orders}
\lean{EnvelopeProperty.GaussLeadingResidueNaturality.verticalOrder_genericField}
\leanok
For every $f\in K$, including zero,
\[
 \operatorname{verticalOrder}_x(J_x^{-1}f)
   =\operatorname{ordFrac}_{D}(f).
\]
In particular the additive integer exponent of a nonzero function is the
SAME at every point $x$, with no ramification or scale factor.
\end{lemma}
\begin{proof}
The existing component theorem makes the SAME $D$ a DVR, by transporting
the existing original DVR $T_x$ through $E_x$. Both are consequently
Noetherian domains of dimension at most one. The existing full-fraction
order transport applied to $E_x$ preserves their normalized native orders.
Its fraction extension is $J_x$: the preceding whole-fraction square
shows equality on all coefficients of $T_x$, and fraction-field uniqueness
extends this to every fraction. This proves the displayed equality.
It includes zero at the level of the order homomorphisms. For nonzero
$f$, injectivity of the integer exponential gives equality of the exact
integer exponents.
\end{proof}

\begin{definition}[One uniformizer and actual global leading functions]
\label{def:gauss-leading-one-original-uniformizer}
\uses{thm:gauss-common-original-component-stalk,lem:gauss-leading-original-vertical-order,lem:gauss-surface-original-residual-vertical-dvrs}
\lean{EnvelopeProperty.GaussLeadingResidueNaturality.exists_globalLeadingFactorization,EnvelopeProperty.GaussLeadingResidueNaturality.exists_parameterUnit}
\leanok
One original point $x_0$ with a Noetherian $S_{x_0}$ and the genuine prime
component equation internally makes $D$ a DVR. Choose ONCE a uniformizer
$\Pi\in D$, independently of all subsequent points and functions.
Every nonzero $f\in K$ has a unique $a_f\in\mathbf Z$ and
$U_f\in D^\times$ such that
\[
 f=\iota_D(U_f)\iota_D(\Pi)^{a_f}.
\]
Define its global leading function $B_f=\lambda(U_f)\in F^\times$.
For every original local component parameter there is an internally
constructed $v_x\in D^\times$ such that
\[
 E_x(\pi_x/1)=\Pi v_x.
\]
Put $V_x=\lambda(v_x)\ne0$. These units and leading functions are
constructed from the actual maps; they are not comparison hypotheses.
\end{definition}
\begin{proof}
Apply the existing component-stalk DVR theorem at $x_0$. Every DVR has an
irreducible element, so choose $\Pi$ once in this SAME $D$.
The DVR fraction factorization theorem writes each nonzero $f$ as a unit
times an integral power of $\iota_D(\Pi)$, allowing negative powers.
The normalized native order of this factorization is $\exp(a_f)$:
the unit has order one and $\Pi$ has additive order one. Hence the
integer exponent is unique. After fixing it, injectivity of $D\to K$
and cancellation of the nonzero power of $\Pi$ make $U_f$ unique.
The original $\lambda$ maps actual units to units of $F$, so $B_f\ne0$.

The original $\pi_x/1$ is irreducible in $T_x$. Its image under the
constructed ring isomorphism $E_x$ is irreducible in $D$. Any two
irreducibles of a DVR are associated. Thus there is an actual unit
$v_x$ with the stated equality; its residue is again nonzero. No
regularity of an ambient closed-point residue field is used.
\end{proof}

\begin{theorem}[The exact local leading-unit transformation]
\label{thm:gauss-leading-local-unit-transformation}
\uses{def:gauss-leading-one-original-uniformizer,lem:gauss-leading-original-localization-squares,lem:gauss-leading-original-vertical-order,thm:gauss-surface-original-fraction-eventual-orders,thm:gauss-common-original-component-stalk}
\lean{EnvelopeProperty.GaussLeadingResidueNaturality.exists_localLeadingFactorization}
\leanok
Fix the ONE $\Pi$. For every nonzero original $f\in K$, the existing
all-fraction producer internally gives $u_{x,f}\in T_x^\times$ and
the SAME integer $a_f$ such that
\[
 J_x^{-1}f=\iota_x(u_{x,f})\pi_x^{a_f},\qquad
 E_x(u_{x,f})=U_fv_x^{-a_f}.
\]
Consequently its unchanged original leading residue is
\[
 r_x(f):=R_x(\operatorname{res}_{T_x}u_{x,f})
       =B_fV_x^{-a_f}\ne0.
\]
Every original nonzero function is covered, including poles in $K$ and
leading functions with negative order at $x$.
\end{theorem}
\begin{proof}
Apply the existing all-fraction producer to $J_x^{-1}f$. It returns a
vertical unit, an integer exponent, its actual native vertical order,
and the genuine leading residue. The preceding whole-order transport
and uniqueness of the global exponent identify the returned exponent
with $a_f$.

Map the returned factorization into $K$ through $J_x$ and use the whole
fraction square. Substitute
$E_x(\pi_x/1)=\Pi v_x$ and the global factorization of $f$.
All powers of $\iota_D(\Pi)$ are nonzero. Cancelling that power gives
\[
 \iota_D(E_xu_{x,f})\iota_D(v_x)^{a_f}=\iota_D(U_f).
\]
Since these are units, multiplication by the inverse power and
injectivity of $D\to K$ yield the claimed whole-unit identity.
Apply the first whole-localization square and the original ring map
$\lambda$; preservation of products and integral powers of units gives
the claimed residue identity. Nonvanishing follows because each of its
factors is nonzero. Thus the local leading residues need not equal
$B_f$; the actual transition unit remains present.
\end{proof}

\begin{lemma}[Positive order of an actual original vertical parameter]
\label{lem:gauss-leading-original-parameter-positive-order}
\uses{def:gauss-leading-one-original-uniformizer,lem:gauss-leading-original-vertical-order,thm:gauss-common-original-component-stalk}
\lean{EnvelopeProperty.GaussLeadingResidueNaturality.exists_verticalParameterFactorization}
\leanok
For an actual $t\in D$ with $t\ne0$ and $t\in\mathfrak m_D$, internally
produce $e\in\mathbf N$, $e>0$, and $U_t\in D^\times$ with
\[
 t=U_t\Pi^e.
\]
For $\tau=\iota_D(t)$ the exact vertical exponent is $e$ and
$B_t=\lambda(U_t)\ne0$. At every original $x$, the existing local
fraction factorization of $J_x^{-1}\tau$ has this SAME exponent and
leading function
$r_x(t)=B_tV_x^{-e}$.
\end{lemma}
\begin{proof}
The DVR factorization of the actual nonzero ring element $t$ gives a
natural exponent $e$ and an actual unit $U_t$. If $e=0$, this equality
makes $t$ a unit, contrary to its membership in the true maximal ideal.
Thus $e>0$. The same factorization in the original fraction field gives
the exact global exponent by uniqueness in the preceding definition.
Apply the local leading-unit transformation to $\tau$. It supplies the
local unit and the displayed residue formula internally.
\end{proof}

Assume now $C$ locally Noetherian, $\dim C\le1$, and $x\ne\eta_C$.
Write $a=a_f$ and
\[
 b_x=\operatorname{ord}_x(r_x(f)),\qquad
 c_x=\operatorname{ord}_x(r_x(t)).
\]
These are integer orders of nonzero functions; neither is assumed
nonnegative. The symbol $c_x$ in this general statement concerns the
local FRACTION representing the given actual generic-stalk parameter.
If that parameter is the germ of an original section near $i(x)$, its
identification with an eventual Gauss-profile coefficient is a separate
consumer of the existing profile theorem.

\begin{theorem}[One genuine whole-curve leading quotient]
\label{thm:gauss-leading-whole-curve-quotient}
\uses{thm:gauss-leading-local-unit-transformation,lem:gauss-leading-original-parameter-positive-order}
\lean{EnvelopeProperty.GaussLeadingResidueNaturality.leadingQuotient_eq_global}
\leanok
The global nonzero function
\[
 H_{f,t}=B_f^e/B_t^a\in F^\times
\]
is defined ONCE before choosing $x$. For every original $x$ with the
genuine local component data, its unchanged local leading functions satisfy
\[
 r_x(f)^e/r_x(t)^a=H_{f,t}.
\]
In particular this SAME function is independent of the local parameter
and second-parameter choices at every $x$, for the ONE fixed $\Pi$.
\end{theorem}
\begin{proof}
Substitute $r_x(f)=B_fV_x^{-a}$ and $r_x(t)=B_tV_x^{-e}$.
Preservation of integer powers and commutativity give transition factors
$V_x^{-ae}$ and $V_x^{-ea}$ in the numerator and denominator. They
are equal and nonzero, so they cancel. This proves the displayed equality
as an equality of actual elements of the SAME $F$.

The right-hand side was defined from the actual global units $U_f,U_t$
before any local point or parameters. The equality holds for EVERY local
choice satisfying the genuine component equations. This is the asserted
independence, proved without assuming equality of local residues.
\end{proof}

\begin{theorem}[The exact corrected original curve order]
\label{thm:gauss-leading-corrected-curve-order}
\uses{thm:gauss-leading-whole-curve-quotient,thm:central-curve-original-leading-order,thm:gauss-leading-local-unit-transformation,lem:gauss-leading-original-parameter-positive-order,thm:fc-curve-proper-support-points,thm:gauss-common-original-component-stalk,thm:gauss-surface-original-fraction-eventual-orders}
\lean{EnvelopeProperty.GaussLeadingResidueNaturality.exists_correctedLeadingOrders}
\leanok
The exponents and orders returned by the unchanged original all-fraction
producer satisfy
\[
 \begin{split}
  e b_x-a c_x&=\operatorname{ord}_x(H_{f,t}),\\
  b_x-\frac ae c_x
   &=\operatorname{ord}_x(B_f)-\frac ae\operatorname{ord}_x(B_t).
 \end{split}
\]
Here the second line is an equality of real numbers and $e>0$ is produced
from the actual parameter. In particular $a/e$ is the SAME vertical
scaling at every point and the transition-unit corrections cancel exactly.
\end{theorem}
\begin{proof}
The existing full-fraction leading-order theorem identifies the actual
native residual order of each produced local unit with the exponential
of the SAME original $C.\operatorname{ord}_x$ of its leading function.
Thus its returned integer is exactly $b_x$, respectively $c_x$, by
injectivity of the integer exponential. This retains the original
local producer's actual units, fractions and coefficients.

All the leading functions are nonzero. The original curve order
homomorphism therefore takes products to sums, integral powers to
integer multiples, and division to subtraction. Apply it to the
preceding whole-field quotient identity. It yields the first displayed
equality. Applying it also to the two local transformation identities,
and putting $d_x=\operatorname{ord}_x(V_x)$, gives
\[
 b_x=\operatorname{ord}_x(B_f)-a d_x,\qquad
 c_x=\operatorname{ord}_x(B_t)-e d_x.
\]
Subtract $a/e$ times the second identity from the first. Since $e>0$,
the terms involving $d_x$ cancel. No sign restriction
on $a$, $b_x$, $c_x$, or $d_x$ was used.

At zero the first two whole-localization squares and the native
whole-field order comparison remain valid, with the zero value of the
order homomorphism. The unit factorizations and these integer formulas
are asserted only for nonzero functions. This avoids identifying the
integer order's value at zero with an exponential residue order.
\end{proof}

The conclusion produces a genuine common function on the whole original
central curve. It does not construct the global surface or its SNC
boundary, alter existing Gauss points, assert a principal-degree balance,
or prove residual nefness or LocalPositivity.
